% !TEX root = tarasov-qasp-zh.tex


\chapter[面镜与透镜的问题你解得好吗？]{面镜与透镜的问题你解得好吗？\label{ch-33}}

\begin{dialogue}

\tchr 我想作几点概括性的说明，它们对于求解涉及透镜与球面镜（凹面镜（concave mirror）与凸面镜（convex mirror））的问题，可能极有用处。这类问题所用的公式可以分为两组。\textbf{第一组}包括把透镜（或面镜）的焦距 $F$、物体到透镜（或面镜）的距离 $d$（物距，object distance）和像到透镜（或面镜）的距离 $f$（像距，image distance）联系起来的公式：
\begin{equation}%
\frac{1}{d}+\frac{1}{f} = \frac{1}{F}
\label{eq-195}
\end{equation}
式中的 $d$、$f$ 和 $F$ 都当作代数量处理，其符号因情形而异。可能的情形只有三种，如下表所列。
\begin{table*}
\begin{small}
\begin{tabular}{p{3.25cm}p{3.25cm}p{4.75cm}}
\toprule%
\multicolumn{2}{l}{会聚透镜与凹面镜} & 发散透镜与凸面镜\\
\midrule%
$d > F$ & $d<F$ & $d > F$\\
$d > 0$、$F > 0$、$f > 0$ & $d>0$、$F>0$、$f<0$ & $d > 0$、$F < 0$、$f < 0$\\
实像 & 虚像 & 虚像\\
\bottomrule%
\end{tabular}
\end{small}
\end{table*}
\pnote{底本表首行作 $d>f$、$d<f$、$d>f$（小写 $f$ 即像距）。但区分这三种情形的条件是物距与焦距的相对位置：会聚元件成实像须 $d>F$、成虚像须 $d<F$，发散元件因 $F<0$ 恒满足 $d>F$；若按像距理解则无法自洽。底本源文件的备选表格正作 $d>F$、$d<F$，当系录入时大小写之误，中文版径改，待考原书。本章另有拼写、标点级小误多处，径改不注。}
于是，$d$ 恒为正；焦距 $F$ 对会聚透镜和凹面镜为正，对发散透镜和凸面镜为负；距离 $f$ 对实像为正、对虚像为负。

\stda 据我理解，利用这张表，可以从普遍公式~\eqref{eq-195} 得出三个用上述各量的算术值（即不带符号的数值）表示的公式：
\begin{equation}
\left. \begin{aligned}
\text{情形 1：}\,\, \frac{1}{d} + \frac{1}{f} &  = \frac{1}{F} \\
\text{情形 2：}\,\, \frac{1}{d}-\frac{1}{f} & = \frac{1}{F} \\
\text{情形 3：}\,\, \frac{1}{d}-\frac{1}{f} & = -\frac{1}{F}
\label{eq-196}
\end{aligned}
\right\}
\end{equation}
\pnote{这套约定与国内教辅处理物像公式的“实正虚负”习惯一致：实像 $f>0$、虚像 $f<0$，会聚元件 $F>0$、发散元件 $F<0$，而 $d$ 恒正。但符号字母与国内不同：国内教材通常以 $u$、$v$、$f$ 分别记物距、像距、焦距，而原书的 $f$ 是像距、焦距记作 $F$，套用公式时须先对准符号。另外，现行课程中透镜成像在初中只通过实验总结规律（人教版八年级上册第五章《透镜及其应用》），物像公式与球面镜公式都不在现行教材正文之内，属竞赛与大学普通物理的内容。}

\tchr 对，正是这样。

\stda 不知怎么，我过去从未留意过透镜与相应球面镜之间的这种类比。

\tchr \textbf{第二组}包括把透镜（或面镜）的焦距与它的其他特征联系起来的公式。对面镜，有简单的关系式
 \begin{equation}%
F = \pm \frac{R}{2}
\label{eq-197}
\end{equation}
式中 $R$ 是面镜的曲率半径（radius of curvature）。正号对应凹面镜（焦距为正），负号对应凸面镜（焦距为负）。对透镜，则有
\begin{equation}%
\frac{1}{F} = (n-1) \left( \frac{1}{R_{1}} + \frac{1}{R_{2}} \right)
\label{eq-198}
\end{equation}
 式中 $n$ 是透镜材料的折射率，$R_{1}$ 和 $R_{2}$ 是透镜两面的曲率半径。若半径 $R$ 属于透镜的凸面，则取正号；若属于凹面，则取负号。容易看出，双凸透镜、平凸透镜（plano-convex lens）和凸凹透镜（convexo-concave lens，会聚弯月形透镜）都是会聚的，因为按公式~\eqref{eq-198}，它们的焦距都是正的。
\pnote{式~\eqref{eq-198} 通常称为磨镜者公式（lensmaker's formula），给出薄透镜（厚度远小于曲率半径）的焦距与两面曲率半径的关系，中学教材不出现，属大学普通物理与竞赛内容。“凸凹透镜必会聚”是有条件的：按式~\eqref{eq-198} 与“凸面取正、凹面取负”的约定，只有凸面曲率半径小于凹面（凸面弯得更厉害）时焦距才为正；反之则得到发散的凹凸透镜（发散弯月形透镜）。原书以“会聚弯月形透镜”标注，指的就是满足这一条件的凸凹透镜。}

\stda 若把透镜放入折射率为 $n_{0}$ 的介质中，公式~\eqref{eq-198} 要作哪些修改？

\tchr 这时应把公式~\eqref{eq-198} 换成
\begin{equation}%
\frac{1}{F} = \left( \frac{n}{n_{0}} - 1\right) \left( \frac{1}{R_{1}} + \frac{1}{R_{2}} \right)
\label{eq-199}
\end{equation}
当周围介质由光疏介质（optically less dense medium，$n_{0}<n$）换成光密介质（optically more dense medium，$n_{0}>n$）时，按公式~\eqref{eq-199}，焦距变号，因而会聚透镜变成发散透镜；反过来，发散透镜也变成会聚透镜。下面我们来解几道具体的题目。
\begin{figure}%[h]%
\centering
%\includegraphics[width=0.5\linewidth]{figs/ch-33/fig-140a}
\input{figs/ch-33/fig-140}
\caption[求此镜的焦距。]{求此镜的焦距。\label{fig-140}}
\end{figure}

\vspace{-20pt}
\hlred{将一块曲率半径为 $R$、折射率为 $n$ 的平凸透镜的凸面镀银，制成一种特殊形式的凹面镜。求此镜的焦距。}
\stda 请允许我来解这道题。我们先让一条光线沿平行于透镜主光轴的方向入射。它经镀银面反射后，又从透镜中射出，并在出射时受到折射。假如不发生这次折射，光线将按公式~\eqref{eq-197} 在距镜面 $R/2$ 处与主光轴相交。由于折射，光线与主光轴的交点离镜面更近一些。把所求的焦距记作 $F$，由\figr{fig-140} 显然有
\begin{equation*}%
\frac{R}{2} \tan \alpha_{1} = F \tan \alpha_{2}
\end{equation*}
由于角 $\alpha_{1}$ 和 $\alpha_{2}$ 都很小，可以应用公式~\eqref{eq-191}。于是
\begin{equation*}%
\frac{R}{2F} = \frac{\tan \alpha_{2}}{\tan \alpha_{1}} \approx \frac{\sin \alpha_{2}}{ \sin \alpha_{1}} = n
\end{equation*}%
由此得
\begin{equation}
F = \frac{R}{2n}
\label{eq-200}
\end{equation}
\stdb 我建议换一种方法来解这道题。大家知道，若把焦距分别为 $F_{1}$ 和 $F_{2}$ 的两个系统组合起来，则新系统的焦距 $F$ 可以按透镜光焦度相加的法则确定，即
\begin{equation}%
\frac{1}{F} = \frac{1}{F_{1}}+\frac{1}{F_{2}}
\label{eq-201}
\end{equation}
\pnote{这里的 power 不是“功率”，而是透镜的光焦度（optical power）$\Phi=1/F$，即焦距的倒数；其单位为 $\mathrm{m}^{-1}$，称屈光度（diopter，D），眼镜的度数即屈光度乘 100。几个彼此贴合的薄光学元件组合时，总光焦度等于各元件光焦度之和，这就是式~\eqref{eq-201}；后文教师把面镜与\emph{两个}透镜的光焦度相加，用的正是这一法则。}
在本题中，按方程~\eqref{eq-198}（其中一个半径为无穷大），透镜的焦距为
\begin{equation*}%
F_{1} = \frac{R}{ (n-1)}
\end{equation*}
而凹面镜的焦距为
\begin{equation*}%
F_{2} = \frac{R}{2}
\end{equation*}
把 $F_{1}$ 和 $F_{2}$ 的表达式代入公式~\eqref{eq-201}，得
\begin{equation}
\frac{1}{F} = \frac{n-1}{R} + \frac{2}{R}
\label{eq-202}
\end{equation}
由此得
\begin{equation}
F = \frac{R}{n+1}
\label{eq-203}
\end{equation}
这说明\smallcaps{学生甲}没有把题做对（他的答案见方程~\eqref{eq-200}）。

\tchr （对\smallcaps{学生乙}）：不对，错的恰恰是你。结果~\eqref{eq-200} 才是正确的。

\stdb 可是，法则~\eqref{eq-201} 在本题情形下难道不对吗？

\tchr 这个法则在本题情形下是正确且适用的。

\stdb 可是，既然法则~\eqref{eq-201} 正确，那么方程~\eqref{eq-202} 也必定正确。

\tchr 你恰恰是在这里弄错了。事实是：光线要两次穿过透镜（一去一回）。因此，应当相加的是面镜和\emph{两个}透镜的光焦度。你不该写方程~\eqref{eq-202}，而应写
\begin{equation*}
\frac{1}{F} = \frac{2(n-1)}{R} + \frac{2}{R}
\end{equation*}
由此得到
\begin{equation*}
\frac{1}{F} = \frac{(2n-2+2)}{R}
\end{equation*}
从而得出
\begin{equation*}
F = \frac{R}{2n}
\end{equation*}
它与方程~\eqref{eq-200} 所得的结果一致。

再看一道题。\hlred{会聚透镜把物体的像放大到 4 倍。若把物体移动 \SI{5}{\centi\meter}，放大率减半。求透镜的焦距。}

\stda 解这类题我总是弄糊涂。我想，必须先画出第一个位置的光路，再画出第二个位置的光路，然后把两者比较一下。

\tchr 我敢说，这道题根本用不着画光路。按公式~\eqref{eq-195}，对给定的位置可以写出
\begin{equation*}%
\frac{1}{F} = \frac{1}{d_{1}} + \frac{1}{f_{1}}
\end{equation*}
 由于第一种情形的放大率（magnification）为 $f_{1}/d_{1} = k_{1}$，我们得到
\begin{equation*}%
\frac{1}{F} = \frac{1}{d_{1}} + \frac{1}{k_{1}d_{1}} = \frac{k_{1}+1}{k_{1}d_{1}}
\end{equation*}
 或
\begin{equation*}%
d_{1} = F \, \frac{k_{1}+1}{k_{1}}
\end{equation*}
同理，对第二个位置可写出
\begin{equation*}%
d_{2} = F \, \frac{k_{2}+1}{k_{2}}
\end{equation*}
于是
\begin{equation}%
d_{2} - d_{1} = F \, \frac{k_{1}-k_{2}}{k_{1}k_{2}}
\label{eq-204}
\end{equation}
按题给条件，$d_{2} - d_{1} = \SI{5}{\centi\metre}, \,\, k_{1}=4 \,\, \text{且} \,\, k_{2} = 2$。把这些值代入方程~\eqref{eq-204}，得 $F = \SI{20}{\centi\metre}$。


\end{dialogue}

\section*{习题}
\label{prob-33}
\addcontentsline{toc}{section}{习题}

\begin{enumerate}[resume*=problems]

\item 一个焦距为 \SI{30}{\centi\metre} 的透镜，成一个缩小到物体大小 2/3 的虚像。它是哪一种透镜（会聚的还是发散的）？物距是多少？若把透镜向远离物体的方向移动 \SI{20}{\centi\metre}，像的大小和像距将是多少？

\item 一个发光点位于曲率半径为 \SI{50}{\centi\metre} 的凹面镜的主光轴上，与镜相距 \SI{15}{\centi\metre}。该点的像在哪里？若把面镜再向远离该点的方向移动 \SI{15}{\centi\metre}，像将怎样？

\item 一个光学系统由一块发散透镜和一块会聚透镜组成（\figr{fig-141}~\drkgry{(a)}，叉号 $\times$ 表示各透镜的焦点）。两透镜的焦距均为 \SI{40}{\centi\metre}。物体位于发散透镜前方 \SI{80}{\centi\metre} 处。试作出该系统所成的物体的像，并计算像的位置。
\begin{figure}[!h]%
\centering
%\includegraphics[width=1.1\linewidth]{
\input{figs/ch-33/fig-141}
\caption[作出系统所成的像并计算其位置。]{作出该系统所成的像，并计算像的位置。\label{fig-141}}
\end{figure}

\item 一个光学系统由三块相同的会聚透镜组成，焦距均为 \SI{30}{\centi\metre}，各透镜的相对位置如\figr{fig-141}~\drkgry{(b)} 所示（叉号 $\times$ 表示各透镜的焦点）。物体与最近的透镜相距 \SI{60}{\centi\metre}。该系统所成的物体的像在哪里？

\item 将一块曲率半径为 \SI{60}{\milli\metre} 的平凸透镜的凸面镀银，制成一个凹面镜。物体位于此镜前方 \SI{25}{\centi\metre} 处。若透镜材料的折射率等于 1.5，求镜到物体的像的距离以及放大率。

\item 将一块曲率半径为 \SI{50}{\centi\metre} 的平凹透镜（plano-concave lens）的凹面镀银，制成一个凸面镜。物体位于此镜前方 \SI{10}{\centi\metre} 处。若透镜材料的折射率等于 1.5，求镜到物体的像的距离以及像的放大率。
\end{enumerate}

\cleardoublepage
% NEWTERMS: concave mirror | 凹面镜 | 首现括注；与凸面镜统称球面镜（ch-32 已定名）
% NEWTERMS: convex mirror | 凸面镜 | 首现括注
% NEWTERMS: radius of curvature | 曲率半径 | 面镜与透镜两面通用；首现括注
% NEWTERMS: plano-convex lens | 平凸透镜 | 首现括注
% NEWTERMS: plano-concave lens | 平凹透镜 | 习题 82 首现括注
% NEWTERMS: convexo-concave lens / (converging) meniscus | 凸凹透镜／（会聚）弯月形透镜 | 首现括注；会聚条件见评注
% NEWTERMS: optically less dense / more dense medium | 光疏介质／光密介质 | 首现括注；人教版选择性必修第一册 4.2《全反射》同名
% NEWTERMS: power (of a lens) | 光焦度 | 焦距的倒数，勿误作“功率”；见评注
% NEWTERMS: diopter | 屈光度 | 仅评注层引入；眼镜度数＝屈光度×100
% NEWTERMS: magnification | 放大率 | 原书符号 $k$，即线放大率 $k=f/d$；正文（第二题解法处）首现括注
% NEWTERMS: object distance / image distance | 物距／像距 | 原书符号 $d$、$f$；国内教材记 $u$、$v$（见评注），正文首现括注
% NEWTERMS: lensmaker's formula | 磨镜者公式 | 仅评注层引入
% NEWTERMS: silver-plated | 镀银 | 普通词，备一致性查考
